\clearpage
% C058：陷阱与纠错，特别说明焦点弦必要条件及水平退化。
\begin{trapbox}
\textbf{\textcolor{CTrap}{陷阱 1：把 $2/p$ 记成 $p/2$}}

\textbf{纠正：}直接代入通径 $x=\frac p2$：两段焦半径都为 $p$，因此倒数和必为 $\frac1p+\frac1p=\frac2p$。\textbf{先用通径复原公式，再做题}，记忆负担最低。

\medskip
\textbf{\textcolor{CTrap}{陷阱 2：斜率设法漏掉通径}}

若写 $y=k(x-\frac p2)$，$k$ 是普通斜率，无法表示竖直直线 $x=\frac p2$，需要单独补通径。\textbf{推荐统一设 $x=my+\frac p2$}，可覆盖所有真实的有限两端点焦点弦，$m=0$ 就是通径。

\medskip
\textbf{\textcolor{CTrap}{陷阱 3：把过焦点的水平直线也当成“弦”}}

$y=0$ 虽然经过焦点 $F(\frac p2,0)$，但联立 $y^2=2px$ 只有 $x=0$ 这\textbf{一个有限交点}，不能取 $A,B$ 两个不同端点。角度法中 $\cos\theta=\pm1$ 分母为 $0$，反映\textbf{有限焦点弦退化}，绝不是出现一个负焦半径。

\medskip
\textbf{\textcolor{CTrap}{陷阱 4：把长度当作有向坐标}}

焦半径 $|AF|,|BF|$ 永远是\textbf{正长度}，不能用有向横坐标位移 $x_A-\frac p2$、$x_B-\frac p2$ 直接替代。右开口抛物线必须用 $|PF|=x_P+\frac p2$；方向负号只体现在 $F$ 两侧的\textbf{投影}中。

\medskip
\textbf{\textcolor{CTrap}{陷阱 5：误以为其他开口方向的定值会变号}}

对于 $y^2=-2px$、$x^2=\pm2py$（统一 $p>0$），通过对称或旋转就回到 $y^2=2px$，\textbf{长度保持不变}，故倒数和仍为 $2/p$。改变的是焦点坐标与焦半径表达式，而不是长度定值。若方程写成 $y^2=-12x$，应取 $p=6$，不是 $p=-6$。

\medskip
\textbf{\textcolor{CTrap}{陷阱 6：从等式倒推出“必过焦点”}}

\textbf{不可逆！}例如 $y^2=4x$，取 $A(\frac14,1)$、$B(4,4)$，虽有 $|AF|=\frac54$、$|BF|=5$，满足 $\frac1{|AF|}+\frac1{|BF|}=1=2/p$，但直线 $AB$ 不过 $F(1,0)$。$x_1x_2=p^2/4$ 也不能单独逆推过焦点，\textbf{需直接检查三点共线}。

\medskip
\textbf{\textcolor{CTrap}{陷阱 7：均值不等式只写下界，忘记取等}}

由 $\frac1r+\frac1s=\frac2p$ 可得 $r+s\ge2p$；\textbf{需要进一步指出} $r=s=p$ 在通径 $x=\frac p2$ 时确实成立，因此最小值才能取到。

\medskip
\textbf{\textcolor{CTrap}{陷阱 8：通径长度与倒数和倍数混乱}}

通径\textbf{整条}长度为 $2p$，每\textbf{半条}长度为 $p$；倒数和 $\frac2p$ 等于整条通径长倒数 $\frac1{2p}$ 的\textbf{四倍}，不是两倍。
\end{trapbox}
